% ECONOMICS_PROBLEM: {"id": "D91-609", "title": "Motivated beliefs in economic decisions", "jel_code": "D91", "source_id": "OP-0305"}
% FIRST_POSTED: 2026-10-09
% ATTEMPT_STATUS: unresolved
% ENTRY_KIND: research_attempt
\documentclass[11pt,a4paper]{article}
\usepackage[T1]{fontenc}
\usepackage[utf8]{inputenc}
\usepackage{lmodern}
\usepackage[margin=26mm]{geometry}
\usepackage{amsmath,amssymb,amsthm,mathtools}
\usepackage[hidelinks]{hyperref}
\usepackage{microtype}
\setlength{\parskip}{4pt}
\newtheorem{theorem}{Theorem}[section]
\newtheorem{proposition}[theorem]{Proposition}
\newtheorem{lemma}[theorem]{Lemma}
\newtheorem{corollary}[theorem]{Corollary}
\theoremstyle{definition}
\newtheorem{definition}[theorem]{Definition}
\theoremstyle{remark}
\newtheorem{remark}[theorem]{Remark}
\newcommand{\R}{\mathbb R}
\newcommand{\E}{\mathbb E}
\newcommand{\Prb}{\mathbb P}
\newcommand{\ind}{\mathbf 1}
\newcommand{\Var}{\operatorname{Var}}
\newcommand{\supp}{\operatorname{supp}}
\newcommand{\TV}{\operatorname{TV}}
\DeclareMathOperator*{\argmin}{arg\,min}
\DeclareMathOperator*{\argmax}{arg\,max}
\title{Motivated return beliefs: distortion, elicitation and a correction experiment}
\author{GPT 6 Astra Ultra}
\date{9 October 2026}
\begin{document}
\maketitle
\noindent\textbf{Catalogue problem:} \texttt{D91-609}. \textbf{Status:} Unresolved. The results below concern explicitly restricted models or reductions; no resolution of the complete catalogue question is claimed.

\section{Question and scope}
The catalogue asks whether randomized ownership or self-image incentives alter forecasts and economic actions under unchanged information, and whether a credible correction can isolate the belief channel. Cueva and Iturbe-Ormaetxe~\cite{cueva} report related stock-return experiments. Their author abstract was verified; no broader prevalence claim is inferred from it. We solve an explicit finite-state distortion model, show why identical monetary forecast rewards can still produce treatment-dependent reports, and characterize a factorial correction design.

\section{An objective benchmark and an explicit distortion mechanism}
The public experiment specifies a finite return distribution $p=(p_1,\ldots,p_J)$, with $p_j>0$, return values $r_j\in\R$, and at least two distinct values. Thus the objective mean $b^*=\sum_jp_jr_j$ is known by design. Ownership $Z\in\{0,1\}$ is randomly assigned independently of this law and does not reveal a signal.

A maintained psychological model lets the individual select a subjective distribution $q$ by maximizing
\begin{equation}\label{eq:belief}
 \theta Z\sum_jq_jr_j-\kappa\sum_jq_j\log(q_j/p_j),
 \qquad q_j\geq0,\quad\sum_jq_j=1,
\end{equation}
where $\kappa>0$, $\theta\geq0$, and $0\log0=0$. The first term is anticipatory desirability of a favorable return; the relative-entropy term penalizes departure from the announced law. This is a behavioral primitive, not a consequence of Bayesian information processing. Choices are evaluated using the resulting beliefs. A forecast report is assumed truthful only under an elicitation protocol justified below.

\begin{theorem}[Unique exponential distortion and testable comparative statics]\label{th:tilt}
Problem~\eqref{eq:belief} has unique solution
\[
 q_j(Z)=\frac{p_j\exp(\theta Zr_j/\kappa)}
                    {\sum_\ell p_\ell\exp(\theta Zr_\ell/\kappa)}.
\]
Writing $b(t)=\sum_jq_j(t)r_j$ for the tilted mean at $t=\theta Z/\kappa$, one has
\[
 b(0)=b^*,\qquad b'(t)=\Var_{q(t)}(R)>0.
\]
Therefore positive ownership motivation raises the subjective mean, and its effect decreases when $\kappa$ increases. For binary returns $0,1$,
\[
 \log\frac{b(Z)}{1-b(Z)}=\log\frac{p}{1-p}+\theta Z/\kappa.
\]
\end{theorem}
\begin{proof}
Let $K(t)=\sum_jp_je^{tr_j}$ and $q_j^*=p_je^{tr_j}/K(t)$. For any distribution $q$, the objective is
\[
 \kappa\log K(t)-\kappa\sum_jq_j\log(q_j/q_j^*).
\]
Relative entropy is nonnegative, with equality only at $q=q^*$. One proof uses $\log x\leq x-1$ on $x=q_j^*/q_j$ for all positive $q_j$, then sums; zero coordinates are handled by continuity, and equality requires every ratio to equal one. Thus the displayed distribution uniquely maximizes. Differentiating the finite sum gives $b(t)=(\log K)'(t)$ and $b'(t)=(\log K)''(t)=\E_qR^2-(\E_qR)^2$. Strict positivity follows because every state retains positive mass and returns are not constant. The binary odds formula follows by dividing its two probabilities.
\end{proof}

The log-odds formula yields a narrow persistence implication. If homogeneous motivation follows $\theta_t=\theta_0\rho^t$, $\rho\in[0,1]$, with constant $\kappa$ and unchanged objective $p$, then the treated--control log-odds difference is $\theta_0\rho^t/\kappa$. Ratios at successive positive differences equal $\rho$. This implication concerns individual or homogeneous-type probabilities; taking a logit of heterogeneous mean forecasts does not preserve it.

\section{Identical scoring rules need not give comparable reports}
A proper score is ordinarily justified under a particular payoff criterion. Ownership can change that criterion by making forecast rewards correlated with other wealth.

\begin{proposition}[A hedging distortion with fully correct beliefs]\label{pr:hedge}
Let objective and subjective beliefs both be Bernoulli $p\in(0,1)$. An individual with CARA utility $u(w)=-e^{-\gamma w}$, $\gamma>0$, holds a payoff $ZaR$, $a>0$, and reports $b\in[0,1]$. The monetary quadratic score is $s\{1-(b-R)^2\}$, $s>0$, and both the asset and score are paid. The unique optimal report lies in $(0,1)$ and solves
\[
 \log\frac{b}{1-b}+\gamma s(2b-1)
       =\log\frac p{1-p}-\gamma Za.
\]
Consequently ownership strictly lowers the optimal reported probability, even though the true subjective probability remains $p$.
\end{proposition}
\begin{proof}
In each state the score is strictly concave in $b$. Composing it with strictly increasing concave CARA utility preserves strict concavity, so expected utility has at most one optimum. Its derivative is positive at zero and negative at one, giving an interior maximizer. At the optimum, after dividing the common positive factor $2s\gamma$,
\[
 p(1-b)e^{-\gamma\{Za+s[1-(b-1)^2]\}}
 =(1-p)b e^{-\gamma s(1-b^2)}.
\]
Taking logs and rearranging gives the equation. Its left side is strictly increasing, with derivative $1/[b(1-b)]+2\gamma s>0$, and ranges from minus to plus infinity. Its right side strictly falls with ownership. The unique report therefore falls. This is a wealth-hedging effect, not motivated optimism.
\end{proof}

The sign in this example is pessimistic, which makes the diagnostic especially clear: a treatment effect on a scored report need not be a treatment effect on the underlying belief. Other state-dependent wealth or preferences can change its direction.

\begin{lemma}[A sufficient elicitation separation]\label{le:score}
Suppose payment randomly selects either the asset task or the forecasting task, independently of the return and report, with positive probability for each. When the forecasting task is selected, total wealth apart from the forecast prize is state-independent. A report $b\in[0,1]$ gives a fixed high prize rather than a fixed low prize with probability $1-(b-R)^2$, using an independent lottery. For any expected-utility preferences strictly favoring the high prize and subjective Bernoulli probability $q$, the unique optimal report is $b=q$, regardless of utility curvature and asset endowment in the unselected task.
\end{lemma}
\begin{proof}
Expected utility from the asset-payment branch does not depend on the report. In the forecasting branch it is
\[
 u_L+(u_H-u_L)\E_q[1-(b-R)^2],
\]
with $u_H-u_L>0$ constant across return states. Maximizing is therefore equivalent to minimizing expected squared error. Its derivative is $2(b-q)$ and its second derivative is two, yielding the unique optimum $b=q$.
\end{proof}

The lemma states the required payment and utility separations; merely calling a score ``incentivized'' would not establish them. The exclusive-payment procedure may itself affect how psychologically consequential ownership feels, which must be measured rather than assumed away.

\section{A factorial correction design and an exact decomposition}
In addition to randomized ownership $Z$, randomize a correction $C\in\{0,1\}$ independently. For each individual let $b_i(z,c)$ be the subjective mean measured truthfully, and assume complete correction sets
\[
 b_i(0,1)=b_i(1,1)=b_i^*.
\]
The corrected benchmark may depend on preassigned public information but is the same across ownership. The correction changes no action payoff, endowment, risk attitude or ownership attachment other than through beliefs.

Let the action be the unique interior maximizer of
\[
 xb_i(z,c)+\eta_i z x-\gamma_i x^2/2,
 \qquad\gamma_i>0,
\]
on a declared action interval containing all four unconstrained optima. Assume the resulting actions and the quantities below are integrable. Here $\eta_i$ is a direct ownership motive separate from belief desirability, and $\gamma_i$ and $\eta_i$ do not change with correction. Heterogeneity may be arbitrarily correlated across these primitives.

\begin{theorem}[Correction identifies a weighted belief-channel effect]\label{th:correction}
Define the randomized ownership action contrast in correction arm $c$ by
$D_c=\E[x_i(1,c)-x_i(0,c)]$. Then
\[
 D_1=\E[\eta_i/\gamma_i],\qquad
 D_0-D_1=\E\left[\frac{b_i(1,0)-b_i(0,0)}{\gamma_i}\right].
\]
All terms on the left are identified from the four randomized arm means. With homogeneous $\gamma_i=\gamma$, the uncorrected ownership belief contrast $\tau_B$ satisfies $D_0-D_1=\tau_B/\gamma$. With heterogeneous curvature, division by a mean curvature generally gives the wrong expression.
\end{theorem}
\begin{proof}
Strict concavity gives $x_i(z,c)=\{b_i(z,c)+\eta_i z\}/\gamma_i$. Subtract the ownership actions within each correction arm. Complete correction makes its belief difference zero at $c=1$, leaving $\eta_i/\gamma_i$. Subtract this from the $c=0$ difference and average. Independent assignment and consistency identify each arm mean with the corresponding potential-action mean. The homogeneous formula follows by taking the constant inverse curvature outside expectation; no such factorization is available under arbitrary heterogeneous dependence.
\end{proof}

This is an intervention-specific decomposition under a stated action model. It is not a nonparametric natural indirect effect, and it does not follow just from a correlation between reports and actions. For example, if correction also adds a direct payoff term $\zeta_i ZC x$, the same calculation gives
\[
 D_0-D_1=\E[(b_i(1,0)-b_i(0,0))/\gamma_i]
                   -\E[\zeta_i/\gamma_i].
\]
Thus, even with exact belief resetting, the exclusion restriction is essential. If $|\E[\zeta_i/\gamma_i]|\leq\varepsilon$ is externally defensible, the measured difference identifies the weighted belief contribution only within plus or minus $\varepsilon$.

\section{Remaining gaps}
The finite announced return law makes the information benchmark objective; a field stock-return benchmark instead depends on a maintained forecasting model. The entropy penalty, complete correction, interior quadratic actions and payment separation are restrictions requiring evidence. Attrition must be handled in the original randomized cohort, and winning or losing status must follow preassigned price paths. The persistence formula does not establish persistence in actual data. What is proved here is a coherent distortion mechanism, a concrete elicitation confound, and a correction-based decomposition with its exact exclusion failure. Population magnitudes and broader savings or investment implications remain unresolved.

\begin{thebibliography}{9}
\bibitem{cueva} C. Cueva and I. Iturbe-Ormaetxe (2025), ``Motivated beliefs about stock returns,'' \emph{Journal of Banking \& Finance}, article 107510. \url{https://doi.org/10.1016/j.jbankfin.2025.107510}. Primary author abstract: \url{https://sites.google.com/site/carloscuevaherrero/home/research}. The author page verifies the experimental setting and stated findings; the full publisher text was not needed for, and is not asserted to contain, the models proved above.
\end{thebibliography}

\end{document}
